\documentclass{zhvt-classic}

% 使用本地源流明体（GenRyuMin2 TW）按文件路径加载，不依赖系统字体名
% 版框（内框／界栏／书口侧栏／鱼尾）改用与版心印章同一印泥棕色：
%   印章图片 yshs.png 实色 = RGB(135,68,52) = HTML 874434
\usepackage{xcolor}
\usepackage{tikz}
% sealbrown 双注册：xcolor 供 TikZ/pgf，L3 供 cls 版框
\definecolor{sealbrown}{HTML}{874434}
\ExplSyntaxOn
\color_set:nnn { sealbrown } { HTML } { 874434 }
\ExplSyntaxOff

\zhvtset{
  grid color       = sealbrown,
  main font        = GenRyuMin2TW-SB,
  main font path   = C:/Users/Administrator/texmf/fonts/opentype/public/GenRyuMin2TW/,
jiazhu font      = GenRyuMin2TW-M,
  jiazhu font path = C:/Users/Administrator/texmf/fonts/opentype/public/GenRyuMin2TW/,
  % 页码框容纳 4 字：正常页「九十七」+前置空 1 格=4 行恰好；
  % 100+ 页「一百零一」共 5 行超出，末字「一/二/三」被框裁掉（2026-10-04 实测），故放宽到 5
  shukou page number chars = 5,
  % 书口书名框高从默认 6 字减到 4 字，把省出的 2 字高拨给 chapter 框，
  % 使 6 字书名+卷名放进 9 字高的 chapter 框后不再压线（2026-10-05）
  shukou title chars = 4,
}

\title{史记}[张耳陈馀列传]
\author{汉·司马迁}
\maker{闲叟漫记·重製}


% ===== 索隐述赞样式（2026-11-08 四次定稿 · 愚大人亲手改定参数）=====
% ⚠ 竖排前提：自定义字体家族必须带 RawFeature={vertical:+vert}；
%   simhei.ttf 加 +vhal 会整体加载失败静默回退，只写 +vert；simfang 同理。
\newCJKfontfamily{\shuzanfont}[Path=C:/Windows/Fonts/,RawFeature={vertical:+vert}]{simhei.ttf}
\newCJKfontfamily{\shuzanbody}[Path=C:/Windows/Fonts/,RawFeature={vertical:+vert}]{simfang.ttf}
% 题头：152x30pt 圆角 6pt 棕底白字；baseline (0,2) 使整框落入界行格子不压版框线；
% 白字节点 x 不补偿、y +8.2pt（像素级墨迹实测 dx/dy≈0，L/R 留白约 20pt）。
\newcommand{\shuzanhead}[1]{%
  \begin{tikzpicture}[baseline={(0,2)}]
    \fill[sealbrown, rounded corners=6pt]
      (-76pt,-15pt) rectangle (76pt,15pt);
    \node[text=white, inner sep=0pt, anchor=center] at (0pt,8.2pt)
      {\shuzanfont\fontsize{23}{23}\selectfont #1};
  \end{tikzpicture}%
}
% ★ 修正书口卷名位置：chapter 框起始位置改回 cls 默认值，避免全局影响。
%   只在每个卷首页通过 \juan 命令里的空命令微调该页位置（安全，不会跨页污染）。
\ExplSyntaxOn
\AtBeginDocument{
  \dim_set:Nn \l__zhvt_shukou_chapter_title_start_dim
    { \l__zhvt_shukou_title_height_dim + \l__zhvt_font_size_dim * 2 }
  \dim_set:Nn \l__zhvt_shukou_page_number_start_dim
    { \l__zhvt_shukou_chapter_title_start_dim + \l__zhvt_shukou_chapter_title_height_dim }
}
\ExplSyntaxOff

\begin{document}

% ★★ 版心（书口）排布覆写：把书名放到鱼尾之下（2026-10-04）
%
%   cls 默认版心自上而下：书名（title 静态框，在鱼尾之上）→ 鱼尾 → 篇名（chapter 动态框）→ 页码
%   本模板改为：          鱼尾（位置不动）→ 书名 → 篇名 → 页码分割线 → 页码
%
%   手法：把 cls 放在 title 静态框里的书名，改放进 chapter 动态框（鱼尾之下）；
%         title 静态框留空。书名取 \l__zhvt_title_str，即 \title{...} 的第一参数，
%         所以换书时只改 \title，不必在此重复写书名。
%
%   ⚠ 三条必须遵守的写法（2026-10-04 实测）：
%     ① 汉字必须用 \hbox{...} 包住。flowfram 会把框内容装进一个宽 14.4mm（= 一个 line_sep
%        = 书口栏宽）的纵向 minipage 并按该宽度折行，而整页靠 PDF /Rotate 90 显示、
%        框内文字本身不变向 —— 不包 \hbox 时「围炉夜话」4 字会被逐字折成 4 行，
%        在版面上就成了横排 4 字并横向溢出版框。
%     ② 竖排下**一个 \hbox 就是一整列**。书名与篇名要同处一列，必须放进同一个 \hbox；
%        分成两个 \hbox 时，后一个会被当成新的一列、甩到隔壁栏去
%        （实测「序」跑到书口栏外侧 14.4mm 的正文列上）。中间的 \quad 在框内是横向空白，
%        显示出来正好是书名与篇名之间的纵向间隔。
%     ③ ★ 内容盒要写成 **与框等宽的定宽盒** `\hbox to \l__zhvt_shukou_chapter_title_height_dim{…\hfil}`。
%        实测：chapter 框是把内容盒沿纵向**居中**摆放的，内容宽度一变，书名就跟着漂：
%        自然宽度的 \hbox（书名 38.4mm）落在 x=157.2mm，加篇名后（57.6mm）落在 x=152.4mm，
%        目录页（无篇名）与正文页的书名就差半格。做成等宽盒后，无论有没有篇名，
%        内容固定从框顶起排，实测两页均落在 x=142.8mm，不再漂。
\ExplSyntaxOn
% 把书口 chapter 框（书名+篇名）起始位置下移 1 字高，避免 4 字卷名顶到鱼尾
% ⚠ cls 的 \__zhvt_calculate_dimensions: 在 \AtBeginDocument 里执行，会重算该 dim；
%   必须把覆写挂到 begin document 之后，否则会被二次覆盖（2026-10-05 实测）。
% (旧方案已移到导言区，直接修改 chapter 框起始位置)
% 卷标题命令 #1=书口篇名(鱼尾) #2=正文卷名（含卷序，卷名内空格用 ~）
%
% ★★ 目录页两项修订（2026-10-05，愚大人指出）★★
%   ① 页码起点统一对齐：
%      原状 —— 目录项写成 \par\hspace{4em}卷名\hspace{4em}页码，页码起点 = 卷名实宽 + 4em。
%      卷名实宽随字数变：「卷一~梁惠王上」6 字 +~、「卷七~离娄上」5 字 +~，
%      差 1 个字位 → 卷七~卷十 4 条的页码整体比其余 10 条高 27.21pt（一个字），看着上下不齐。
%      改法 —— 卷名补全角空格（宽 = 1 字 = 27.215pt）补足到 **7 字符位**再写进目录，
%      14 条页码起点遂全部落在同一条竖线上（起点对齐；末字随页码字数自然参差）。
%      ⚠ 只对目录里的副本 \l_tmpc_tl 补，正文卷名仍用 \l_tmpa_tl 原样 —— 正文像素不受影响。
%   ② 卷名整体上移半字位：前导 \hspace 由 4em 改 3.5em。
%      实测原状目录块在版框内偏下（天头留白 139.6pt / 地脚留白 110.7pt），
%      上移 0.5em = 13.6pt 后上下基本均分，视觉居中。
\NewDocumentCommand { \juan } { mm } {
	\clearpage
	\tl_set:Nn \l__zhvt_chapter_title_tl { #1 }
	\tl_set:Nn \l_tmpa_tl { #2 }
	\tl_set_eq:NN \l_tmpc_tl \l_tmpa_tl
	\int_step_inline:nn { 7 - \str_count:n { #2 } }
	  { \tl_put_right:Nn \l_tmpc_tl { 　 } }
	\tl_set:Nx \l_tmpb_tl { \zhnum { page } }
	\tl_put_left:Nn \l_tmpb_tl { \hspace { 4em } }
	\tl_put_left:Nn \l_tmpb_tl { \l_tmpc_tl }
	\tl_put_left:Nn \l_tmpb_tl { \par \hspace { 3.5em } }
	\addtocontents { toc } { \l_tmpb_tl }
	\skip_horizontal:n { \l__zhvt_font_size_dim * 4 }
	\l_tmpa_tl
	\par
	\thispagestyle { empty }
}
\cs_set_protected:Npn \__zhvt_set_shukou_content: {
	% ★ 鱼尾上方留白（title 静态框）嵌入闲叟钤印：
	%   整页 /Rotate 90，故此处把图用 \rotatebox{90} 转正为竖章；
	%   尺寸卡在版心条幅内（书内约 7mm 宽 × 33.8mm 高，留余量不压鱼尾）。
	\setstaticcontents * { title } {
		\vbox_to_ht:nn { \l__zhvt_baseline_skip_dim } {
			\vskip 14.2pt
			\hbox_to_wd:nn { \l__zhvt_shukou_title_height_dim } {
				\hfil
				\rotatebox [ origin = c ] { 90 } { \includegraphics [ width = 7mm , height = 33.8mm ] { yshs.png } }
				\hfil
			}
		}
	}
% ★ 版心（书口）统一改法（2026-11-16 愚大人定）：18pt＋框内居中＋书口文字同版框色
%   ① 原 27.2pt 时 7 字长篇名（如「廉颇蔺相如列传」）末字溢出 chapter 框、压住页码；
%      改 \fontsize{18pt}{18pt}，前后 \hfil 令字串在框内居中（不随篇名字数漂移）。
%   ② \vbox{\vskip 8.2pt} 定量下移，抵消缩字号后栏宽方向 8.2pt 上偏（像素实测 y 中心 543.3 = 原版）。
%      ⚠ \raisebox / \vspace* 在 flowfram 框内无效；换字号须同步换补偿：16pt→10.1、19pt→7.4。
%   ③ chapter 与 page~number 均置 \color{sealbrown}（HTML 874434），与版框／界行同色；正文仍黑。
%   ⚠ 页码字位、钤印、鱼尾、界行、正文一律未动；五部合订本同此口径。
	\setdynamiccontents * { chapter }
	{ \color{sealbrown} \vbox { \vskip 8.2pt \hbox to \l__zhvt_shukou_chapter_title_height_dim
		{ \hfil \fontsize{18pt}{18pt}\selectfont \l__zhvt_title_str \quad \l__zhvt_chapter_title_tl \hfil } } }
	\setdynamiccontents * { page~number } { \color{sealbrown} \quad \zhnum { page } }
}
\__zhvt_set_shukou_content:
\ExplSyntaxOff

\maketitle

% 单篇，无多卷目录，不渲染目录页

\juan{张耳陈馀列传}{卷八十九~张耳陈馀列传第二十九}
\hspace*{2em}张耳者。大梁人也。其少时。及魏公子毋忌为客。张耳尝亡命游外黄。外黄富人女甚美。嫁庸奴。亡其夫。去抵父客。父客素知张耳。乃谓女曰。﹁必欲求贤夫。从张耳。﹂女听。乃卒为请决。嫁之张耳。张耳是时脱身游。女家厚奉给张耳。张耳以故致千里客。乃宦魏为外黄令。名由此益贤。陈馀者。亦大梁人也。好儒术。数游赵苦陉。富人公乘氏以其女妻之。亦知陈馀非庸人也。馀年少。父事张耳。两人相与为刎颈交。\\
\hspace*{2em}秦之灭大梁也。张耳家外黄。高祖为布衣时。尝数从张耳游。客数月。秦灭魏数岁。已闻此两人魏之名士也。购求有得张耳千金。陈馀五百金。张耳。陈馀乃变名姓。俱之陈。为里监门以自食。两人相对。里吏尝有过笞陈馀。陈馀欲起。张耳蹑之。使受笞。吏去。张耳乃引陈馀之桑下而数之曰。﹁始吾与公言何如。今见小辱而欲死一吏乎。﹂陈馀然之。秦诏书购求两人。两人亦反用门者以令里中。\\
\hspace*{2em}陈涉起蕲。至入陈。兵数万。张耳。陈馀上谒陈涉。涉及左右生平数闻张耳。陈馀贤。未尝见。见即大喜。\\
\hspace*{2em}陈中豪杰父老乃说陈涉曰。﹁将军身被坚执锐。率士卒以诛暴秦。复立楚社稷。存亡继绝。功德宜为王。且夫监临天下诸将。不为王不可。原将军立为楚王也。﹂陈涉问此两人。两人对曰。﹁夫秦为无道。破人国家。灭人社稷。绝人後世。罢百姓之力。尽百姓之财。将军瞋目张胆。出万死不顾一生之计。为天下除残也。今始至陈而王之。示天下私。原将军毋王。急引兵而西。遣人立六国後。自为树党。为秦益敌也。敌多则力分。与众则兵彊。如此野无交兵。县无守城。诛暴秦。据咸阳以令诸侯。诸侯亡而得立。以德服之。如此则帝业成矣。今独王陈。恐天下解也。﹂陈涉不听。遂立为王。\\
\hspace*{2em}陈馀乃复说陈王曰。﹁大王举梁。楚而西。务在入关。未及收河北也。臣尝游赵。知其豪桀及地形。原请奇兵北略赵地。﹂於是陈王以故所善陈人武臣为将军。邵骚为护军。以张耳。陈馀为左右校尉。予卒三千人。北略赵地。\\
\hspace*{2em}武臣等从白马渡河。至诸县。说其豪桀曰。﹁秦为乱政虐刑以残贼天下。数十年矣。北有长城之役。南有五岭之戍。外内骚动。百姓罢敝。头会箕敛。以供军费。财匮力尽。民不聊生。重之以苛法峻刑。使天下父子不相安。陈王奋臂为天下倡始。王楚之地。方二千里。莫不响应。家自为怒。人自为斗。各报其怨而攻其雠。县杀其令丞。郡杀其守尉。今已张大楚。王陈。使吴广。周文将卒百万西击秦。於此时而不成封侯之业者。非人豪也。诸君试相与计之。夫天下同心而苦秦久矣。因天下之力而攻无道之君。报父兄之怨而成割地有土之业。此士之一时也。﹂豪桀皆然其言。乃行收兵。得数万人。号武臣为武信君。下赵十城。馀皆城守。莫肯下。\\
\hspace*{2em}乃引兵东北击范阳。范阳人蒯通说范阳令曰。﹁窃闻公之将死。故吊。虽然。贺公得通而生。﹂范阳令曰。﹁何以吊之。﹂对曰。﹁秦法重。足下为范阳令十年矣。杀人之父。孤人之子。断人之足。黥人之首。不可胜数。然而慈父孝子莫敢倳刃公之腹中者。畏秦法耳。今天下大乱。秦法不施。然则慈父孝子且倳刃公之腹中以成其名。此臣之所以吊公也。今诸侯畔秦矣。武信君兵且至。而君坚守范阳。少年皆争杀君。下武信君。君急遣臣见武信君。可转祸为福。在今矣。﹂\\
\hspace*{2em}范阳令乃使蒯通见武信君曰。﹁足下必将战胜然後略地。攻得然後下城。臣窃以为过矣。诚听臣之计。可不攻而降城。不战而略地。传檄而千里定。可乎。﹂武信君曰。﹁何谓也。﹂蒯通曰。﹁今范阳令宜整顿其士卒以守战者也。怯而畏死。贪而重富贵。故欲先天下降。畏君以为秦所置吏。诛杀如前十城也。然今范阳少年亦方杀其令。自以城距君。君何不赍臣侯印。拜范阳令。范阳令则以城下君。少年亦不敢杀其令。令范阳令乘硃轮华毂。使驱驰燕。赵郊。燕。赵郊见之。皆曰此范阳令。先下者也。即喜矣。燕。赵城可毋战而降也。此臣之所谓传檄而千里定者也。﹂武信君从其计。因使蒯通赐范阳令侯印。赵地闻之。不战以城下者三十馀城。\\
\hspace*{2em}至邯郸。张耳。陈馀闻周章军入关。至戏卻。又闻诸将为陈王徇地。多以谗毁得罪诛。怨陈王不用其筴不以为将而以为校尉。乃说武臣曰。﹁陈王起蕲。至陈而王。非必立六国後。将军今以三千人下赵数十城。独介居河北。不王无以填之。且陈王听谗。还报。恐不脱於祸。又不如立其兄弟。不。即立赵後。将军毋失时。时间不容息。﹂武臣乃听之。遂立为赵王。以陈馀为大将军。张耳为右丞相。邵骚为左丞相。\\
\hspace*{2em}使人报陈王。陈王大怒。欲尽族武臣等家。而发兵击赵。陈王相国房君谏曰。﹁秦未亡而诛武臣等家。此又生一秦也。不如因而贺之。使急引兵西击秦。﹂陈王然之。从其计。徙系武臣等家宫中。封张耳子敖为成都君。\\
\hspace*{2em}陈王使使者贺赵。令趣发兵西入关。张耳。陈馀说武臣曰。﹁王王赵。非楚意。特以计贺王。楚已灭秦。必加兵於赵。原王毋西兵。北徇燕。代。南收河内以自广。赵南据大河。北有燕。代。楚虽胜秦。必不敢制赵。﹂赵王以为然。因不西兵。而使韩广略燕。李良略常山。张黡略上党。\\
\hspace*{2em}韩广至燕。燕人因立广为燕王。赵王乃与张耳。陈馀北略地燕界。赵王间出。为燕军所得。燕将囚之。欲与分赵地半。乃归王。使者往。燕辄杀之以求地。张耳。陈馀患之。有厮养卒谢其舍中曰。﹁吾为公说燕。与赵王载归。﹂舍中皆笑曰。﹁使者往十馀辈。辄死。若何以能得王。﹂乃走燕壁。燕将见之。问燕将曰。﹁知臣何欲。﹂燕将曰。﹁若欲得赵王耳。﹂曰。﹁君知张耳。陈馀何如人也。﹂燕将曰。﹁贤人也。﹂曰。﹁知其志何欲。﹂曰。﹁欲得其王耳。﹂赵养卒乃笑曰。﹁君未知此两人所欲也。夫武臣。张耳。陈馀杖马箠下赵数十城。此亦各欲南面而王。岂欲为卿相终己邪。夫臣与主岂可同日而道哉。顾其势初定。未敢参分而王。且以少长先立武臣为王。以持赵心。今赵地已服。此两人亦欲分赵而王。时未可耳。今君乃囚赵王。此两人名为求赵王。实欲燕杀之。此两人分赵自立。夫以一赵尚易燕。况以两贤王左提右挈。而责杀王之罪。灭燕易矣。﹂燕将以为然。乃归赵王。养卒为御而归。\\
\hspace*{2em}李良已定常山。还报。赵王复使良略太原。至石邑。秦兵塞井陉。未能前。秦将诈称二世使人遗李良书。不封。曰。﹁良尝事我得显幸。良诚能反赵为秦。赦良罪。贵良。﹂良得书。疑不信。乃还之邯郸。益请兵。未至。道逢赵王姊出饮。从百馀骑。李良望见。以为王。伏谒道旁。王姊醉。不知其将。使骑谢李良。李良素贵。起。惭其从官。从官有一人曰。﹁天下畔秦。能者先立。且赵王素出将军下。今女兒乃不为将军下车。请追杀之。﹂李良已得秦书。固欲反赵。未决。因此怒。遣人追杀王姊道中。乃遂将其兵袭邯郸。邯郸不知。竟杀武臣。邵骚。赵人多为张耳。陈馀耳目者。以故得脱出。收其兵。得数万人。客有说张耳曰。﹁两君羁旅。而欲附赵。难。独立赵後。扶以义。可就功。﹂乃求得赵歇。立为赵王。居信都。李良进兵击陈馀。陈馀败李良。李良走归章邯。\\
\hspace*{2em}章邯引兵至邯郸。皆徙其民河内。夷其城郭。张耳与赵王歇走入钜鹿城。王离围之。陈馀北收常山兵。得数万人。军钜鹿北。章邯军钜鹿南棘原。筑甬道属河。饷王离。王离兵食多。急攻钜鹿。钜鹿城中食尽兵少。张耳数使人召前陈馀。陈馀自度兵少。不敌秦。不敢前。数月。张耳大怒。怨陈馀。使张黡。陈泽往让陈馀曰。﹁始吾与公为刎颈交。今王与耳旦暮且死。而公拥兵数万。不肯相救。安在其相为死。苟必信。胡不赴秦军俱死。且有十一二相全。﹂陈馀曰。﹁吾度前终不能救赵。徒尽亡军。且馀所以不俱死。欲为赵王。张君报秦。今必俱死。如以肉委饿虎。何益。﹂张黡。陈泽曰。﹁事已急。要以俱死立信。安知後虑。﹂陈馀曰。﹁吾死顾以为无益。必如公言。﹂乃使五千人令张黡。陈泽先尝秦军。至皆没。\\
\hspace*{2em}当是时。燕。齐。楚闻赵急。皆来救。张敖亦北收代兵。得万馀人。来。皆壁馀旁。未敢击秦。项羽兵数绝章邯甬道。王离军乏食。项羽悉引兵渡河。遂破章邯。章邯引兵解。诸侯军乃敢击围钜鹿秦军。遂虏王离。涉间自杀。卒存钜鹿者。楚力也。\\
\hspace*{2em}於是赵王歇。张耳乃得出钜鹿。谢诸侯。张耳与陈馀相见。责让陈馀以不肯救赵。及问张黡。陈泽所在。陈馀怒曰。﹁张黡。陈泽以必死责臣。臣使将五千人先尝秦军。皆没不出。﹂张耳不信。以为杀之。数问陈馀。陈馀怒曰。﹁不意君之望臣深也。岂以臣为重去将哉。﹂乃脱解印绶。推予张耳。张耳亦愕不受。陈馀起如厕。客有说张耳曰。﹁臣闻﹃天与不取。反受其咎﹄。今陈将军与君印。君不受。反天不祥。急取之。﹂张耳乃佩其印。收其麾下。而陈馀还。亦望张耳不让。遂趋出。张耳遂收其兵。陈馀独与麾下所善数百人之河上泽中渔猎。由此陈馀。张耳遂有卻。\\
\hspace*{2em}赵王歇复居信都。张耳从项羽诸侯入关。汉元年二月。项羽立诸侯王。张耳雅游。人多为之言。项羽亦素数闻张耳贤。乃分赵立张耳为常山王。治信都。信都更名襄国。\\
\hspace*{2em}陈馀客多说项羽曰。﹁陈馀。张耳一体有功於赵。﹂项羽以陈馀不从入关。闻其在南皮。即以南皮旁三县以封之。而徙赵王歇王代。\\
\hspace*{2em}张耳之国。陈馀愈益怒。曰。﹁张耳与馀功等也。今张耳王。馀独侯。此项羽不平。﹂及齐王田荣畔楚。陈馀乃使夏说说田荣曰。﹁项羽为天下宰不平。尽王诸将善地。徙故王王恶地。今赵王乃居代。原王假臣兵。请以南皮为扞蔽。﹂田荣欲树党於赵以反楚。乃遣兵从陈馀。陈馀因悉三县兵袭常山王张耳。张耳败走。念诸侯无可归者。曰。﹁汉王与我有旧故。而项羽又彊。立我。我欲之楚。﹂甘公曰。﹁汉王之入关。五星聚东井。东井者。秦分也。先至必霸。楚虽彊。後必属汉。﹂故耳走汉。汉王亦还定三秦。方围章邯废丘。张耳谒汉王。汉王厚遇之。\\
\hspace*{2em}陈馀已败张耳。皆复收赵地。迎赵王於代。复为赵王。赵王德陈馀。立以为代王。陈馀为赵王弱。国初定。不之国。留傅赵王。而使夏说以相国守代。\\
\hspace*{2em}汉二年。东击楚。使使告赵。欲与俱。陈馀曰。﹁汉杀张耳乃从。﹂於是汉王求人类张耳者斩之。持其头遗陈馀。陈馀乃遣兵助汉。汉之败於彭城西。陈馀亦复觉张耳不死。即背汉。\\
\hspace*{2em}汉三年。韩信已定魏地。遣张耳与韩信击破赵井陉。斩陈馀泜水上。追杀赵王歇襄国。汉立张耳为赵王。汉五年。张耳薨。谥为景王。子敖嗣立为赵王。高祖长女鲁元公主为赵王敖后。\\
\hspace*{2em}汉七年。高祖从平城过赵。赵王朝夕袒韝蔽。自上食。礼甚卑。有子婿礼。高祖箕踞詈。甚慢易之。赵相贯高。赵午等年六十馀。故张耳客也。生平为气。乃怒曰。﹁吾王孱王也。﹂说王曰。﹁夫天下豪桀并起。能者先立。今王事高祖甚恭。而高祖无礼。请为王杀之。﹂张敖齧其指出血。曰。﹁君何言之误。且先人亡国。赖高祖得复国。德流子孙。秋豪皆高祖力也。原君无复出口。﹂贯高。赵午等十馀人皆相谓曰。﹁乃吾等非也。吾王长者。不倍德。且吾等义不辱。今怨高祖辱我王。故欲杀之。何乃汙王为乎。令事成归王。事败独身坐耳。﹂\\
\hspace*{2em}汉八年。上从东垣还。过赵。贯高等乃壁人柏人。要之置厕。上过欲宿。心动问曰。﹁县名为何。﹂曰。﹁柏人。﹂﹁柏人者。迫於人也。﹂不宿而去。\\
\hspace*{2em}汉九年。贯高怨家知其谋。乃上变告之。於是上皆并逮捕赵王。贯高等。十馀人皆争自刭。贯高独怒骂曰。﹁谁令公为之。今王实无谋。而并捕王。公等皆死。谁白王不反者。﹂乃轞车胶致。与王诣长安。治张敖之罪。上乃诏赵群臣宾客有敢从王皆族。贯高与客孟舒等十馀人。皆自髡钳。为王家奴。从来。贯高至。对狱。曰。﹁独吾属为之。王实不知。﹂吏治榜笞数千。刺剟。身无可击者。终不复言。吕后数言张王以鲁元公主故。不宜有此。上怒曰。﹁使张敖据天下。岂少而女乎。﹂不听。廷尉以贯高事辞闻。上曰。﹁壮士。谁知者。以私问之。﹂中大夫泄公曰。﹁臣之邑子。素知之。此固赵国立名义不侵为然诺者也。﹂上使泄公持节问之箯舆前。仰视曰。﹁泄公邪。﹂泄公劳苦如生平驩。与语。问张王果有计谋不。高曰。﹁人情宁不各爱其父母妻子乎。今吾三族皆以论死。岂以王易吾亲哉。顾为王实不反。独吾等为之。﹂具道本指所以为者王不知状。於是泄公入。具以报。上乃赦赵王。\\
\hspace*{2em}上贤贯高为人能立然诺。使泄公具告之。曰。﹁张王已出。﹂因赦贯高。贯高喜曰。﹁吾王审出乎。﹂泄公曰。﹁然。﹂泄公曰。﹁上多足下。故赦足下。﹂贯高曰。﹁所以不死一身无馀者。白张王不反也。今王已出。吾责已塞。死不恨矣。且人臣有篡杀之名。何面目复事上哉。纵上不杀我。我不愧於心乎。﹂乃仰绝肮。遂死。当此之时。名闻天下。\\
\hspace*{2em}张敖已出。以尚鲁元公主故。封为宣平侯。於是上贤张王诸客。以钳奴从张王入关。无不为诸侯相。郡守者。及孝惠。高后。文帝。孝景时。张王客子孙皆得为二千石。\\
\hspace*{2em}张敖。高后六年薨。子偃为鲁元王。以母吕后女故。吕后封为鲁元王。元王弱。兄弟少。乃封张敖他姬子二人。寿为乐昌侯。侈为信都侯。高后崩。诸吕无道。大臣诛之。而废鲁元王及乐昌侯。信诸侯。孝文帝即位。复封故鲁元王偃为南宫侯。续张氏。\\
\hspace*{2em}太史公曰。张耳。陈馀。世传所称贤者。其宾客厮役。莫非天下俊桀。所居国无不取卿相者。然张耳。陈馀始居约时。相然信以死。岂顾问哉。及据国争权。卒相灭亡。何乡者相慕用之诚。後相倍之戾也。岂非以势利交哉。名誉虽高。宾客虽盛。所由殆与大伯。延陵季子异矣。\\

% ===== 索隐述赞（司马贞《史记索隐》）· 统一样式 =====
% 题头「索隐述赞曰」五字并入反白框；正文每四句一列，定宽盒 \hfil 居列中；
% \fontsize{24}{41} 列距对齐全书界行（121+41k）；\clearpage 使述赞独占一页。
\clearpage
\shuzanhead{索隐述赞曰}\linebreak
{\shuzanbody\fontsize{24}{41}\selectfont
\hbox to \hsize{\hfil 张耳。陈馀。天下豪俊。忘年羁旅。\hfil}\linebreak
\hbox to \hsize{\hfil 刎颈相信。耳围钜鹿。馀兵不进。张既望深。\hfil}\linebreak
\hbox to \hsize{\hfil 陈乃去印。势利倾夺。隙末成衅。\hfil}%
\par}
\end{document}
